\documentclass[14pt]{extarticle}
% ---------- PACKAGES ----------
\usepackage{amsmath, amssymb}
\usepackage{geometry}
\usepackage{setspace}
\usepackage{titlesec}
\usepackage{xcolor}
\usepackage{helvet}
\usepackage{enumitem}
\usepackage{lmodern}
\makeatletter
\IfFileExists{../common-things/reflectionpage.tex}{%
  \input{../common-things/reflectionpage.tex}%
}{%
  \IfFileExists{common-things/reflectionpage.tex}{%
    \input{common-things/reflectionpage.tex}%
  }{%
    \PackageError{\jobname}{Missing reflectionpage.tex file}{%
      Place reflectionpage.tex inside a common-things/ directory next to this unit\@.%
    }%
  }%
}
\makeatother
% ---------- PAGE SETUP ----------
\geometry{margin=1in}
\setstretch{1.5}
\setlength{\parindent}{0pt}
\setlength{\parskip}{0.75em}
\setlist{itemsep=0.75\baselineskip, topsep=0.5\baselineskip}
\titleformat{\section}{\normalfont\Large\bfseries}{\thesection}{1em}{}
\titleformat{\subsection}{\normalfont\large\bfseries}{\thesubsection}{1em}{}
\renewcommand{\familydefault}{\sfdefault}
\renewcommand{\emph}[1]{\textbf{#1}}

\begin{document}
\raggedright
\pagenumbering{gobble}

\begin{center}
    \LARGE \textbf{Unit 9: Multi-Step and Mixed Problems} \\[6pt]
    \Large \textbf{Topic 2: Combining Systems, Functions, and Graphs}
\end{center}

\vspace{1em}

\section*{Concept Summary}

This topic brings together three core algebra domains: \textbf{systems of equations}, \textbf{functions and their properties}, and \textbf{coordinate-plane graphs}. On the SAT, harder problems frequently mix these areas. You might need to solve a system of equations to find a parameter that appears in a function rule, or interpret the intersection of two graphs as the solution to a system.

\textbf{Systems meet functions:} A common pattern is two functions \(f(x)\) and \(g(x)\) whose graphs intersect. Setting \(f(x) = g(x)\) gives you an equation to solve, and the solution is the \(x\)-coordinate of the intersection point. If the functions are linear, you get a system of two linear equations. If one or both are quadratic, solving \(f(x) = g(x)\) becomes a quadratic equation, and the number of solutions tells you whether the graphs intersect zero, one, or two times.

\textbf{Functions meet graphs:} The SAT often describes a function with an equation and then asks about features of its graph: intercepts, vertex, slope, or end behavior. Going the other direction, a problem might give you a graph and ask you to write or identify the equation. Key connections to keep in mind:

\begin{itemize}
  \item The \(y\)-intercept of \(f(x)\) is \(f(0)\).
  \item The \(x\)-intercepts (zeros) are the solutions of \(f(x) = 0\).
  \item The vertex of a quadratic \(f(x) = ax^2 + bx + c\) is at \(x = -\dfrac{b}{2a}\).
  \item A system of one linear and one quadratic equation can have 0, 1, or 2 solutions. Use the discriminant of the resulting quadratic to determine which.
\end{itemize}

\textbf{Systems with parameters:} Some problems include a constant (like \(k\)) in one of the equations and ask for the value of \(k\) that makes the system have exactly one solution, no solution, or infinitely many solutions. For two linear equations, parallel lines (no solution) have equal slopes but different intercepts. For a line and a parabola, exactly one intersection means the discriminant of the combined equation is zero.

\textbf{Function composition and transformation:} If \(f(x) = 2x + 1\) and \(g(x) = x^2\), then \(f(g(3)) = f(9) = 19\). The SAT tests this at a straightforward level. Transformations (shifts, reflections, stretches) connect a function's equation to how its graph moves: \(f(x - 3)\) shifts right by 3, \(f(x) + 5\) shifts up by 5, and \(-f(x)\) reflects over the \(x\)-axis.

\section*{Core Skills}
\begin{itemize}
  \item Find the intersection point(s) of two functions by setting them equal and solving.
  \item Use the discriminant to determine the number of intersections between a line and a parabola.
  \item Evaluate composite functions and apply function notation in multi-step problems.
  \item Connect features of a function's equation (slope, vertex, intercepts) to features of its graph.
  \item Find the value of a parameter in a system that produces a specified number of solutions.
\end{itemize}

\section*{Example 1: Intersection of Two Linear Functions}

Let \(f(x) = 3x + 2\) and \(g(x) = -x + 10\). At what point do the graphs of \(f\) and \(g\) intersect?

\textbf{Step 1:} Set \(f(x) = g(x)\).
\[
3x + 2 = -x + 10 \implies 4x = 8 \implies x = 2
\]

\textbf{Step 2:} Find the \(y\)-coordinate.
\[
f(2) = 3(2) + 2 = 8
\]

The intersection point is \(\boxed{(2, 8)}\).

\section*{Example 2: Line and Parabola System}

For what value of \(k\) does the line \(y = 2x + k\) intersect the parabola \(y = x^2\) at exactly one point?

\textbf{Step 1:} Set the equations equal.
\[
x^2 = 2x + k \implies x^2 - 2x - k = 0
\]

\textbf{Step 2:} For exactly one intersection, the discriminant equals zero.
\[
b^2 - 4ac = (-2)^2 - 4(1)(-k) = 4 + 4k = 0 \implies k = -1
\]

The value of \(k\) is \(\boxed{-1}\).

\section*{Example 3: Composite Functions}

If \(f(x) = x^2 - 1\) and \(g(x) = 2x + 3\), what is \(f(g(2))\)?

\textbf{Step 1:} Evaluate the inner function first.
\[
g(2) = 2(2) + 3 = 7
\]

\textbf{Step 2:} Evaluate the outer function at that result.
\[
f(7) = 7^2 - 1 = 49 - 1 = 48
\]

\(\boxed{48}\)

\section*{Example 4: System with No Solution}

For what value of \(a\) does the system below have no solution?
\[
y = 3x + 5 \qquad y = ax - 1
\]

\textbf{Step 1:} Two lines have no solution when they are parallel: same slope, different \(y\)-intercepts. The first line has slope 3 and \(y\)-intercept 5.

\textbf{Step 2:} For the lines to be parallel, \(a = 3\). Since the \(y\)-intercepts (\(5\) and \(-1\)) are already different, \(a = 3\) gives no solution.

The value of \(a\) is \(\boxed{3}\).

\section*{Example 5: Graph Features from a Quadratic}

The function \(f(x) = -2(x - 3)^2 + 8\) is graphed in the \(xy\)-plane. What is the maximum value of \(f\), and at what value of \(x\) does it occur?

\textbf{Step 1:} The function is in vertex form \(a(x - h)^2 + k\). The vertex is \((3, 8)\). Since \(a = -2 < 0\), the parabola opens downward, so the vertex is a maximum.

The maximum value is \(\boxed{8}\), occurring at \(x = \boxed{3}\).

\section*{Example 6: System Solved by Substitution with a Quadratic}

Solve the system: \(y = x + 1\) and \(y = x^2 - 3x + 5\).

\textbf{Step 1:} Set the expressions equal.
\[
x + 1 = x^2 - 3x + 5 \implies 0 = x^2 - 4x + 4 = (x - 2)^2
\]

\textbf{Step 2:} Solve.
\[
x - 2 = 0 \implies x = 2
\]
\[
y = 2 + 1 = 3
\]

The system has one solution: \(\boxed{(2, 3)}\).

\section*{Key Takeaways}
\begin{itemize}
  \item The intersection of two graphs is found by setting the functions equal. The algebraic method always matches the geometric picture.
  \item The discriminant \(b^2 - 4ac\) is your tool for counting intersections between a line and a parabola. Positive gives two intersections, zero gives one (tangent), and negative gives none.
  \item For parallel lines (no solution), match slopes; for identical lines (infinitely many solutions), match both slope and intercept.
  \item When evaluating composite functions, always work inside-out: evaluate the innermost function first, then use that result as the input for the outer function.
\end{itemize}

\newpage

\section*{Practice Questions: Combining Systems, Functions, and Graphs}

\subsection*{Part A: Intersections of Functions}
\begin{enumerate}
  \item Let \(f(x) = 2x - 1\) and \(g(x) = -x + 8\). What is the \(x\)-coordinate of the intersection point of the graphs of \(f\) and \(g\)?
  \item Let \(f(x) = x^2\) and \(g(x) = 4x - 4\). Find all \(x\)-values where \(f(x) = g(x)\).
  \item The line \(y = 5\) intersects the parabola \(y = x^2 - 4\). What are the \(x\)-coordinates of the two intersection points?
  \item Let \(f(x) = x^2 + 2x\) and \(g(x) = 3x + 6\). What is the sum of the \(x\)-coordinates of the intersection points?
  \item Let \(h(x) = \dfrac{12}{x}\) and \(k(x) = x + 1\). Find all \(x\)-values where \(h(x) = k(x)\). (Assume \(x \neq 0\).)
\end{enumerate}

\subsection*{Part B: Discriminant and Number of Solutions}
\begin{enumerate}
  \item The line \(y = 4x + c\) is tangent to the parabola \(y = x^2\). What is the value of \(c\)?
  \item For what values of \(m\) does the line \(y = mx\) intersect the parabola \(y = x^2 + 1\) at exactly two points?
  \item The system \(y = -x + k\) and \(y = x^2 + 2x + 3\) has no solution. Find all values of \(k\) for which this is true. Express your answer as an inequality.
  \item For what value of \(b\) does the system \(y = x^2 + bx + 5\) and \(y = 2x + 1\) have exactly one solution?
  \item The graphs of \(y = x^2 + 3\) and \(y = 2x + k\) intersect at exactly one point. What is the \(x\)-coordinate of the intersection?
\end{enumerate}

\subsection*{Part C: Composite and Evaluated Functions}
\begin{enumerate}
  \item If \(f(x) = 3x + 2\) and \(g(x) = x^2\), what is \(g(f(1))\)?
  \item If \(f(x) = x - 4\) and \(g(x) = 2x^2 + 1\), what is \(f(g(3))\)?
  \item If \(f(x) = 5x - 3\), for what value of \(x\) is \(f(f(x)) = 32\)?
  \item Let \(f(x) = x^2 + 1\). What is \(f(a + 1) - f(a)\)?
  \item If \(f(x) = \dfrac{x + 1}{x - 1}\) for \(x \neq 1\), what is \(f(f(2))\)?
\end{enumerate}

\subsection*{Part D: Graph Features and Transformations}
\begin{enumerate}
  \item The graph of \(f(x) = (x - 4)(x + 2)\) crosses the \(x\)-axis at two points. What is the sum of the \(x\)-coordinates of these points?
  \item The vertex of \(g(x) = x^2 - 6x + 5\) is at the point \((h, k)\). What is the value of \(k\)?
  \item The graph of \(y = f(x)\) is shifted 3 units to the right and 2 units down to produce \(y = g(x)\). If \(f(x) = x^2\), what is \(g(x)\)?
  \item The function \(f(x) = -3(x + 1)^2 + 12\) has a maximum value. At what \(x\)-value does this maximum occur, and what is the maximum value?
  \item The graph of \(f(x) = |x|\) is reflected over the \(x\)-axis and then shifted up by 5 units. Write the equation of the resulting function and find the \(y\)-value when \(x = 3\).
\end{enumerate}

\subsection*{Part E: SAT-Style Applications}
\begin{enumerate}
  \item A company's revenue, in thousands of dollars, is modeled by \(R(x) = -2x^2 + 20x\), and its cost is modeled by \(C(x) = 4x + 10\), where \(x\) is the number of units (in thousands) produced. For what values of \(x\) does revenue equal cost?
  \item A ball is thrown upward from a 48-foot platform with an initial velocity of 32 feet per second. Its height after \(t\) seconds is given by \(h(t) = -16t^2 + 32t + 48\). At the same time, an elevator rises from the ground at a constant rate, with its height given by \(e(t) = 16t\). At what time \(t > 0\) are the ball and elevator at the same height?
  \item The function \(f(x) = ax^2 + bx + c\) passes through the points \((0, 5)\), \((1, 2)\), and \((2, 3)\). What is the value of \(a + b + c\)?
  \item The line \(y = mx - 3\) intersects the curve \(y = x^2 + x + 1\) at exactly one point. Find the two possible values of \(m\) and their sum.
  \item Let \(f(x) = x^2 - 4x + 3\). The graph of \(f\) is a parabola. A horizontal line \(y = c\) intersects the parabola at exactly one point. What is the value of \(c\)?
\end{enumerate}

\newpage

\section*{Answer Key and Solutions: Combining Systems, Functions, and Graphs}

\subsection*{Part A Solutions: Intersections of Functions}
\begin{enumerate}
  \item Set \(2x - 1 = -x + 8\):
  \[
  3x = 9 \implies x = 3
  \]

  \(\boxed{3}\)

  \item Set \(x^2 = 4x - 4\):
  \[
  x^2 - 4x + 4 = 0 \implies (x - 2)^2 = 0 \implies x = 2
  \]

  \(\boxed{2}\)

  \item Set \(x^2 - 4 = 5\):
  \[
  x^2 = 9 \implies x = 3 \text{ or } x = -3
  \]

  \(\boxed{3 \text{ and } -3}\)

  \item Set \(x^2 + 2x = 3x + 6\):
  \[
  x^2 - x - 6 = 0 \implies (x - 3)(x + 2) = 0 \implies x = 3 \text{ or } x = -2
  \]
  The sum is \(3 + (-2) = 1\).

  \(\boxed{1}\)

  \item Set \(\dfrac{12}{x} = x + 1\). Multiply by \(x\):
  \[
  12 = x^2 + x \implies x^2 + x - 12 = 0 \implies (x + 4)(x - 3) = 0
  \]
  \(x = -4\) or \(x = 3\).

  \(\boxed{-4 \text{ and } 3}\)
\end{enumerate}

\subsection*{Part B Solutions: Discriminant and Number of Solutions}
\begin{enumerate}
  \item Set \(x^2 = 4x + c\), giving \(x^2 - 4x - c = 0\). For tangency, discriminant \(= 0\):
  \[
  16 + 4c = 0 \implies c = -4
  \]

  \(\boxed{-4}\)

  \item Set \(mx = x^2 + 1\), giving \(x^2 - mx + 1 = 0\). For two solutions, discriminant \(> 0\):
  \[
  m^2 - 4 > 0 \implies m^2 > 4 \implies m > 2 \text{ or } m < -2
  \]

  \(\boxed{m > 2 \text{ or } m < -2}\)

  \item Set \(-x + k = x^2 + 2x + 3\), giving \(x^2 + 3x + (3 - k) = 0\). For no solution, discriminant \(< 0\):
  \[
  9 - 4(3 - k) < 0 \implies 9 - 12 + 4k < 0 \implies 4k < 3 \implies k < \frac{3}{4}
  \]

  \(\boxed{k < \dfrac{3}{4}}\)

  \item Set \(x^2 + bx + 5 = 2x + 1\), giving \(x^2 + (b - 2)x + 4 = 0\). For one solution, discriminant \(= 0\):
  \[
  (b - 2)^2 - 16 = 0 \implies (b - 2)^2 = 16 \implies b - 2 = \pm 4
  \]
  \(b = 6\) or \(b = -2\). Both are valid; the problem asks for ``the value,'' so the answer includes both.

  \(\boxed{6 \text{ or } -2}\)

  \item Set \(x^2 + 3 = 2x + k\), giving \(x^2 - 2x + (3 - k) = 0\). For one intersection, discriminant \(= 0\):
  \[
  4 - 4(3 - k) = 0 \implies 4 - 12 + 4k = 0 \implies k = 2
  \]
  Now solve \(x^2 - 2x + 1 = 0 \implies (x - 1)^2 = 0 \implies x = 1\).

  \(\boxed{1}\)
\end{enumerate}

\subsection*{Part C Solutions: Composite and Evaluated Functions}
\begin{enumerate}
  \item \(f(1) = 3(1) + 2 = 5\). Then \(g(5) = 5^2 = 25\).

  \(\boxed{25}\)

  \item \(g(3) = 2(9) + 1 = 19\). Then \(f(19) = 19 - 4 = 15\).

  \(\boxed{15}\)

  \item First, \(f(x) = 5x - 3\). Then \(f(f(x)) = 5(5x - 3) - 3 = 25x - 15 - 3 = 25x - 18\).
  \[
  25x - 18 = 32 \implies 25x = 50 \implies x = 2
  \]

  \(\boxed{2}\)

  \item \(f(a + 1) = (a+1)^2 + 1 = a^2 + 2a + 2\). \(f(a) = a^2 + 1\).
  \[
  f(a+1) - f(a) = (a^2 + 2a + 2) - (a^2 + 1) = 2a + 1
  \]

  \(\boxed{2a + 1}\)

  \item \(f(2) = \dfrac{2 + 1}{2 - 1} = 3\). Then \(f(3) = \dfrac{3 + 1}{3 - 1} = \dfrac{4}{2} = 2\).

  \(\boxed{2}\)
\end{enumerate}

\subsection*{Part D Solutions: Graph Features and Transformations}
\begin{enumerate}
  \item The \(x\)-intercepts are \(x = 4\) and \(x = -2\). Their sum is \(4 + (-2) = 2\).

  \(\boxed{2}\)

  \item The \(x\)-coordinate of the vertex is \(x = -\dfrac{-6}{2(1)} = 3\). Then \(k = g(3) = 9 - 18 + 5 = -4\).

  \(\boxed{-4}\)

  \item Shifting right 3: replace \(x\) with \(x - 3\). Shifting down 2: subtract 2.
  \[
  g(x) = (x - 3)^2 - 2
  \]

  \(\boxed{(x-3)^2 - 2}\)

  \item The vertex is at \(x = -1\) (from the form \(-3(x + 1)^2 + 12\)). The maximum value is 12.

  Maximum at \(x = \boxed{-1}\), maximum value \(= \boxed{12}\).

  \item Reflecting over the \(x\)-axis: \(y = -|x|\). Shifting up 5: \(y = -|x| + 5\). At \(x = 3\): \(y = -3 + 5 = 2\).

  \(\boxed{2}\)
\end{enumerate}

\subsection*{Part E Solutions: SAT-Style Applications}
\begin{enumerate}
  \item Set \(R(x) = C(x)\):
  \[
  -2x^2 + 20x = 4x + 10 \implies -2x^2 + 16x - 10 = 0 \implies x^2 - 8x + 5 = 0
  \]
  By the quadratic formula:
  \[
  x = \frac{8 \pm \sqrt{64 - 20}}{2} = \frac{8 \pm \sqrt{44}}{2} = \frac{8 \pm 2\sqrt{11}}{2} = 4 \pm \sqrt{11}
  \]

  \(\boxed{4 + \sqrt{11} \text{ and } 4 - \sqrt{11}}\)

  \item Set \(h(t) = e(t)\):
  \[
  -16t^2 + 32t + 48 = 16t \implies -16t^2 + 16t + 48 = 0
  \]
  Divide by \(-16\):
  \[
  t^2 - t - 3 = 0 \implies t = \frac{1 \pm \sqrt{1 + 12}}{2} = \frac{1 \pm \sqrt{13}}{2}
  \]
  Since \(t > 0\), we take the positive root: \(t = \dfrac{1 + \sqrt{13}}{2}\).

  \(\boxed{\dfrac{1 + \sqrt{13}}{2}}\)

  \item From \((0, 5)\): \(c = 5\). From \((1, 2)\): \(a + b + 5 = 2 \implies a + b = -3\). So \(a + b + c = -3 + 5 = 2\).

  (We can verify with \((2, 3)\): \(4a + 2b + 5 = 3 \implies 4a + 2b = -2\). Combined with \(a + b = -3\), we get \(a = 2, b = -5\), and \(f(2) = 8 - 10 + 5 = 3\). This checks out.)

  \(\boxed{2}\)

  \item Set \(mx - 3 = x^2 + x + 1\), giving \(x^2 + (1 - m)x + 4 = 0\). For exactly one intersection, the discriminant equals zero:
  \[
  (1 - m)^2 - 4(1)(4) = 0 \implies (1 - m)^2 = 16 \implies 1 - m = \pm 4
  \]
  So \(m = -3\) or \(m = 5\). Their sum is \(-3 + 5 = 2\).

  \(\boxed{2}\)

  \item The minimum value of \(f(x) = x^2 - 4x + 3\) occurs at the vertex. The \(x\)-coordinate is \(x = \dfrac{4}{2} = 2\). Then \(f(2) = 4 - 8 + 3 = -1\). A horizontal line is tangent to the parabola at the vertex, so \(c = -1\).

  \(\boxed{-1}\)
\end{enumerate}

\ReflectionPage{Unit 9, Topic 2}{https://www.scotchegg.co}

\end{document}
